\section{측정 데이터와 스트림 형식}
\label{sec:data}
\subsection{물리 차원, 단위와 명목 측정량}
물리 차원을 SI 기본량의 지수 벡터로 표현한다.
\begin{equation}
 d=(d_M,d_L,d_T,d_I,d_\Theta,d_N,d_J)\in\mathbb Z^7.
\end{equation}
예를 들어 전압의 차원은 $(1,2,-3,-1,0,0,0)$, 주파수의 차원은
$(0,0,-1,0,0,0,0)$이다. 곱셈은 차원 벡터를 더하고 나눗셈은 차원 벡터를 뺀다.
단위 레코드 $u=(d,s,o)$는 기준 단위와의 관계 $x_{\mathrm{SI}}=sx+o$를 제공한다.
전기생리 전압과 농도 변화는 각 연산이 사용하는 선형 단위 표현을 선언한다.

채널 $c$의 기술자를 다음과 같이 정의한다.
\begin{equation}
 \chi_c=(\mathrm{id}_c,q_c,d_c,u_c,r_c,g_c,\mathrm{role}_c),
 \qquad \mathcal C=(\chi_1,\ldots,\chi_C).
 \label{eq:channels}
\end{equation}
$q_c$는 명목 측정량, $r_c$는 기준, $g_c$는 위치ㆍ기하 정보이다.
전극 전압의 식별자는 측정 전극과 기준 전극 또는 기준 연산을 포함한다.
광학 측정의 식별자는 광원, 검출기, 파장과 측정 종류를 포함한다.
기하 정보는 좌표값과 함께 좌표계 식별자, 길이 단위, 장치 배치를 보관한다.

\begin{table}[tb]
\centering\small
\caption{채널 역할별 의미와 형식 정보}
\begin{tabularx}{\linewidth}{@{}P{0.19\linewidth}YY@{}}
\toprule
역할 & 측정량과 단위 & 식별 정보 \\
\midrule
전기생리 & 전극 전압, V 또는 $\mu$V & 전극, 기준, 위치, 증폭 설정 \\
광학 원시값 & 광강도에 대응하는 장치 단위 & 광원, 검출기, 파장, 스케일 \\
광학 파생값 & 광밀도 또는 농도 변화 & 측정 종류, 단위, 광경로 모형 \\
운동 센서 & 가속도, 각속도 & 축, 좌표계, 센서 위치 \\
디지털 사건 & 상태 변화, trigger word & 입출력 선, 극성, 관측 시계 \\
스파이크 사건 & 발생 시각과 선택적 파형 & 전극ㆍunit 식별자, 검출기 버전 \\
\bottomrule
\end{tabularx}
\label{tab:channel_roles}
\end{table}

\subsection{표본 발생 규칙과 수치 표현}
\begin{definition}[스트림 형식]
스트림의 형식을 다음 다섯 성분으로 정의한다.
\begin{equation}
 \tau=\Stream[T_{\mathrm{val}},\mathcal C,\kappa,\mathcal R,\mathcal Q].
 \label{eq:streamtype}
\end{equation}
$T_{\mathrm{val}}$은 값의 형식, $\kappa$는 시계 영역,
$\mathcal R$은 표본 발생 규칙, $\mathcal Q$는 품질 어휘이다.
수치 기본형 $D$는 정수ㆍ부동소수점ㆍ패킹 표현을 정한다.
값 형식은 $D$, $\operatorname{Vec}[D,C]$, $\operatorname{Epoch}[D,C,N]$,
$\operatorname{SPD}[D,C]$, 사건의 합 형식으로 구성한다.
벡터와 epoch의 채널 축은 $\mathcal C$의 순서를 사용한다.
\begin{equation}
 \mathcal R=\operatorname{Uniform}(f,\phi)
 \mid\operatorname{Events},\qquad f\in\mathbb Q_{>0}.
\end{equation}
균일 스트림은 값 레코드의 명목 발생률 $f$와 위상 $\phi$를 사용한다.
epoch 값은 내부 표본 격자와 별도로 epoch 발생률을 가진다.
이벤트 스트림은 각 레코드의 시간 정보를 사용한다.
\end{definition}

획득 블록의 값 배열은 $X\in D^{C\times N}$으로 표현한다. 첫 번째 축은 채널,
두 번째 축은 표본이다. 어댑터는 입력 버퍼의 stride, endianness와 채널 순서를
이 형식에 대응시킨다. 배열의 저장 표현과 채널의 물리 단위는 각자의 기술자를 유지한다.
BrainFlow의 행렬 경로는 EEG, 보조 센서, 사건, timestamp 행을 제공하며,
각 preset이 자신의 표본률과 시간 정보를 가진다\cite{brainflow_data}.
이 경로의 어댑터는 행 역할을 분해하여 식~\eqref{eq:streamtype}의 스트림들을 생성한다.

\subsection{측정 시각, 가용 시각과 세대}
\begin{definition}[시간 레코드]
한 표본 또는 사건을 다음 튜플로 나타낸다.
\begin{equation}
 r_i=(\mathrm{id}_i,x_i,K_i,I_i^{\mathrm{meas}},a_i,
       e_i^{\mathrm{schema}},e_i^{\mathrm{clock}},q_i,\rho_i).
 \label{eq:record}
\end{equation}
$I_i^{\mathrm{meas}}$는 대표 측정 시각의 구간이고, $a_i$는 이름이 지정된 관측 경계에서
해당 레코드를 이용할 수 있게 된 시각이다. $\rho_i$는 데이터ㆍ매개변수ㆍ변환의 계보를 나타낸다.
\end{definition}

시계 변환이 점 추정 $\widehat t_i$와 오차 반경 $\epsilon_i$를 제공하면
\begin{equation}
 I_i^{\mathrm{meas}}=[\widehat t_i-\epsilon_i,\widehat t_i+\epsilon_i]
\end{equation}
로 표현한다. 획득 패킷 도착, SDK 반환, bridge 송신, DSL 수신은 각각 관측 경계 이름을 갖는다.
연산 결과는 대표 시각을 유지하면서 자신의 실제 가용 시각을 기록한다.
데이터의 전체 의존 시간 지지는 원천 조상의 측정 구간으로부터 계산한다.
\begin{equation}
 \operatorname{Supp}(y)=
 \bigcup_{r\in\operatorname{SourceAnc}(y)}I_r^{\mathrm{meas}}.
 \label{eq:support}
\end{equation}

표본 인덱스의 출처는 다음 합 형식으로 보관한다.
\begin{equation}
 K_i=\operatorname{Observed}(k_i)
 \mid\operatorname{Reconstructed}(k_i,\nu_i)
 \mid\operatorname{Unresolved}.
\end{equation}
$\nu_i$는 카운터, 표본률, 보정 구간 등 재구성에 사용한 근거이다.
수신기는 acquisition epoch에서 표본 순서의 연속성을 관리하고,
schema epoch에서 채널ㆍ스케일ㆍ표현의 연속성을 관리한다.
clock epoch는 장치 시계의 재시작과 변환 유효 구간을 식별한다.

\begin{definition}[가용 접두 스트림]
실행 환경과 안정적인 레코드 순서를 고정하고, 시각 $u$에 가용한 부분을 다음과 같이 정의한다.
\begin{equation}
 \cut_u(S)=(r_i\in S\mid a_i\le u).
 \label{eq:cut}
\end{equation}
동일한 가용 시각의 레코드는 원천 식별자와 순번으로 정렬한다.
타이머, 설정 변경, 모델 활성화도 이름이 지정된 입력 사건으로 포함한다.
\end{definition}

\subsection{품질 상태와 계보 그래프}
품질 상태는 $q_i=(v_i,o_i,F_i)$로 구성한다. $v_i$는 채널별 유효성 마스크,
$o_i$는 표본 출처, $F_i$는 사건 플래그 집합이다. 출처는
$\mathsf{Measured}$, $\mathsf{Interpolated}$, $\mathsf{Recovered}$ 등을 사용한다.
플래그는 gap, saturation, decoder resynchronization, epoch transition을 표현한다.
Cortex의 EEG 레코드는 보간 여부를 제공하고\cite{cortex_sample}, OxySoft는
연결 손실 자료의 복구 경로를 제공한다\cite{oxysoft}. 어댑터는 이 정보를
품질 상태와 가용 시각에 반영한다.

계보는 방향성 비순환 그래프로 저장한다. 원천 노드에는 장치 프로필과 원본 범위를,
변환 노드에는 연산자ㆍ매개변수ㆍ버전과 입력 노드 식별자를 기록한다.
모델 노드는 학습 자료, 전처리 통계와 선택 절차를 연결한다.
구현은 정규화된 레코드의 내용 해시를 식별자로 사용하여 반복 참조를 공유한다.
수치 연산과 메타데이터 변환이 동일한 계보 경로를 따라가므로,
출력에서 원래 측정과 활성 모델까지의 관계를 복원할 수 있다.
